\begin{table}[htbp]
\centering
\caption{Standardized Effect Sizes}
\label{tab:sde}
\begin{threeparttable}
\small
\begin{adjustbox}{max width=\textwidth}
\begin{tabular}{lcccccc}
\toprule
Outcome & $\hat{\beta}$ & SE & SD($Y$) & SDE & SE(SDE) & Classification \\
\midrule
\multicolumn{7}{l}{\textit{Panel A: Pooled}} \\
Oil heating share & -0.1474 & 0.2456 & 6.61 & -0.2209 & 0.368 & Large negative \\
Heat pump share & 0.1755 & 0.1696 & 2.08 & 0.8373 & 0.8092 & Large positive \\
Gas heating share & 0.3259 & 0.0917 & 9.51 & 0.3394 & 0.0955 & Large positive \\
\midrule
\multicolumn{7}{l}{\textit{Panel B: By Gas Infrastructure}} \\
Gas share (high gas infra.) & 0.3863 & 0.0948 & 4.78 & 0.8306 & 0.2038 & Large positive \\
Gas share (low gas infra.) & 0.358 & 0.1915 & 4.88 & 0.6761 & 0.3615 & Large positive \\
\bottomrule
\end{tabular}
\end{adjustbox}
\begin{tablenotes}[flushleft]
\footnotesize
\item \textit{Notes:} \textbf{Country:} Switzerland. \textbf{Research question:} Does the federal CO2 levy on heating fuels cause cantonal-level fuel switching from oil to gas or heat pumps? \textbf{Policy mechanism:} Switzerland's CO2 levy imposes a per-ton charge on fossil heating fuels (oil and gas), increasing from CHF 12/ton in 2008 to CHF 120/ton in 2022, funded by a legislated ratchet that triggers increases when national emission reduction targets are missed. One-third of revenue funds building renovation subsidies. \textbf{Outcome definition:} Share of dwellings (\%) using each energy source for primary heating, from the BFS GWS register. \textbf{Treatment:} Continuous; canton's 2000 Census oil-heating share (proportion) multiplied by the federal levy rate (CHF/ton CO2). \textbf{Data:} BFS Geb\"aude- und Wohnungsstatistik (GWS), 2000 and 2021--2024, 26 cantons, 127 canton-year observations. \textbf{Method:} Two-way fixed effects (canton + year), standard errors clustered at canton level. \textbf{Sample:} All 26 Swiss cantons; oil heating share in 2000 ranges from 41\% (Basel-Stadt) to 74\% (Jura). SDE $= \hat{\beta} \times \text{SD}(X) / \text{SD}(Y)$ where SD($X$) is the standard deviation of treatment intensity among post-treatment observations and SD($Y$) is the pre-treatment (2000) standard deviation of the outcome. Classification refers to magnitude, not statistical significance: Large ($|$SDE$| > 0.15$), Moderate ($0.05$--$0.15$), Small ($0.005$--$0.05$), Null ($< 0.005$).
\end{tablenotes}
\end{threeparttable}
\end{table}

